\subsection{Definiciones, Acrónimos y Abreviaciones}

En la Tabla \ref{tab:acronimos} se definen los términos técnicos, acrónimos y abreviaciones utilizados a lo largo de este documento para garantizar una interpretación uniforme entre todos los lectores.

\begin{longtable}{|p{3.2cm}|p{11.5cm}|}
\caption{Definiciones, acrónimos y abreviaciones} \label{tab:acronimos} \\
\hline
\textbf{Término / Sigla} & \textbf{Definición} \\ \hline
\endfirsthead

\multicolumn{2}{c}{\tablename\ \thetable\ -- \textit{Continuación}} \\ \hline
\textbf{Término / Sigla} & \textbf{Definición} \\ \hline
\endhead

\hline \multicolumn{2}{r}{\textit{Continúa en la siguiente página}} \\
\endfoot

\hline
\endlastfoot

\textbf{API} & \textit{Application Programming Interface}. Interfaz de programación de aplicaciones utilizada para consultar datos meteorológicos externos. \\ \hline
\textbf{Borde activo} & Franja o perímetro del incendio donde continúa activa la combustión y el avance del frente de llama. \\ \hline
\textbf{CC BY 4.0} & Licencia Creative Commons Atribución 4.0 Internacional bajo la cual se distribuyen los datos meteorológicos de Open-Meteo. \\ \hline
\textbf{CFD} & \textit{Computational Fluid Dynamics}. Dinámica de fluidos computacional empleada para modelar el comportamiento físico del fuego y el humo. \\ \hline
\textbf{CU} & Caso de Uso definido en el Plan de Administración del Proyecto (PMP). \\ \hline
\textbf{DAD} & \textit{Disciplined Agile Delivery}. Marco de trabajo ágil híbrido adoptado para el desarrollo iterativo del proyecto. \\ \hline
\textbf{FDS} & \textit{Fire Dynamics Simulator}. Modelo computacional de dinámica de fluidos desarrollado por el NIST para simular incendios y transporte de humo. \\ \hline
\textbf{HPC-ZINE} & Infraestructura de cómputo de alto desempeño utilizada durante la fase de desarrollo para precalcular \textcolor{black}{el dataset de simulaciones}. \\ \hline
\textbf{HTTP / HTTPS} & Protocolo de transferencia de hipertexto (y su variante segura) utilizado para realizar consultas GET a la API meteorológica. \\ \hline
\textbf{JSON} & \textit{JavaScript Object Notation}. Formato ligero de intercambio de datos en el que se reciben las respuestas de Open-Meteo y se estructuran estimaciones. \\ \hline
\textbf{MD} & Elemento del Modelo del Dominio del sistema. \\ \hline
\textbf{NIST} & \textit{National Institute of Standards and Technology}. Instituto responsable del desarrollo de FDS y Smokeview. \\ \hline
\textbf{PMP} & \textit{Project Management Plan}. Plan de Administración del Proyecto donde se especifican casos de uso, cronograma y gestión. \\ \hline
\textbf{Punto de ignición} & Coordenada geográfica de referencia donde se origina el fuego en un escenario de simulación o consulta. \\ \hline
\textbf{SMV / Smokeview} & Herramienta de visualización desarrollada por el NIST para representar gráficamente los resultados tridimensionales generados por FDS. \\ \hline
\textbf{SQL / SQLite} & Lenguaje de consulta estructurada y motor de base de datos relacional embebido utilizado para almacenar el catálogo de escenarios. \\ \hline
\textbf{SRS} & \textit{Software Requirements Specification}. Especificación de Requerimientos de Software. \\ \hline
\textbf{SUS} & \textit{System Usability Scale}. Escala estandarizada utilizada para evaluar la facilidad de uso de la interfaz con usuarios reales. \\ \hline
\textbf{VFP} & Visión y Formulación del Proyecto (Propuesta de trabajo de grado). \\ \hline
\end{longtable}